\documentclass[11pt]{article}
\usepackage[margin=1in]{geometry}
\usepackage{xcolor}
\usepackage{hyperref}
\usepackage{booktabs}
\usepackage{parskip}

\definecolor{revcomment}{rgb}{0.2,0.2,0.6}
\definecolor{response}{rgb}{0,0.4,0}

\newcommand{\reviewer}[1]{\noindent\textbf{\large Reviewer #1}\medskip\hrule\medskip}
\newcommand{\rcomment}[1]{\noindent\textcolor{revcomment}{\textit{Reviewer comment:}}~#1\par\medskip}
\newcommand{\rresponse}[1]{\noindent\textcolor{response}{\textbf{Response:}}~#1\par\medskip\medskip}

\title{Response to Reviewers\\
\large Manuscript ID: Access-2025-56953\\
\normalsize ``Position-Aware Versus Semantic Chunking for Content Generation: Small Gains, Big Trade-Offs''}
\author{Gourav Singla\\\small Independent Researcher}
\date{}

\begin{document}
\maketitle

I thank the Associate Editor (Dr.~Shadi Alawneh) and both reviewers for their detailed and constructive feedback on the original submission. Their comments significantly strengthened this revised manuscript. The major revisions are summarized first, followed by point-by-point responses to each reviewer.

\section*{Summary of Major Revisions}

In response to reviewer feedback, the manuscript has been substantially expanded along every dimension flagged as a limitation:

\begin{enumerate}
    \item \textbf{Multiple generation models.} The original paper used only GPT-4o. The revision adds Claude Sonnet~4 and Gemini~2.0~Flash, totaling three independent LLMs from three different providers. Method rankings are consistent across all three (see the new ``Cross-Generator Consistency'' subsection in Results).
    \item \textbf{Multiple downstream tasks.} The original paper evaluated only slide generation. The revision adds factual question answering as a second downstream task with 10 questions per document and an objective correctness metric (1{,}741 evaluations, approximately 17{,}000 question judgments). Q\&A reveals dramatically larger chunking-method differences than slides (see the new ``Question Answering: Truncation Dominates'' subsection in Results).
    \item \textbf{Multiple LLM judges.} The original paper used only Gemini~2.0~Flash as judge. The revision adds GPT-4o-mini and Claude~Sonnet~4 as additional independent judges, with cross-judge agreement analysis (new ``Cross-Judge Agreement'' subsection in Results). Three judges from three model families produce identical rankings.
    \item \textbf{Human evaluation.} 50 blind document pairs were rated by an expert annotator, validating the LLM-as-judge findings (new ``Human Evaluation Validates LLM Judges'' subsection in Results).
    \item \textbf{Expanded dataset.} The corpus grew from 101 research papers to 218 diverse documents, including non-research material (government reports, NGO whitepapers, industry outlooks, technical manuals, educational materials) spanning more than 20 domains.
    \item \textbf{Theoretical framework.} A new ``Theoretical Framework'' section formalizes the information-density model, derives the W-curve from lost-in-the-middle theory, and predicts the conditions under which position-aware chunking should outperform truncation. These theoretical predictions are then empirically tested.
    \item \textbf{Ablation studies.} Chunking was re-run at four budgets (1{,}000, 2{,}000, 3{,}000, and 5{,}000 words) on all 218 documents (new ``Ablation: Chunk Budget Sensitivity'' subsection).
    \item \textbf{W-curve sensitivity.} Position-aware chunking was re-run with 11 alternative parameter configurations (new ``W-Curve Parameter Sensitivity'' subsection) showing the method is robust to parameter perturbation.
    \item \textbf{Algorithm and prompt details.} Algorithm~1 provides full pseudocode for the position-aware chunking selection procedure; the slide evaluation prompt is included as a typeset block in the Methodology section.
    \item \textbf{Speed-quality trade-off.} The roughly 3{,}300$\times$ latency advantage of truncation over semantic chunking is now stated in the abstract, expanded into a dedicated ``Computational Cost (Latency)'' subsection in Results, and discussed as a practical implication in the new ``Practical Implications: The Speed-Quality Trade-off'' subsection of the Discussion.
    \item \textbf{Novel finding.} Length-stratified analysis (new ``Position-Aware Chunking Has a Niche: Medium-Length Documents'' subsection) reveals that position-aware chunking outperforms all alternatives for medium-length documents (5K--10K words) on slide generation, identifying its empirical niche.
    \item \textbf{Practitioner guidelines} (new ``When to Use What: A Practitioner Guide'' subsection in the Discussion) translate findings into task-aware recommendations.
\end{enumerate}

The revised manuscript reflects approximately 6{,}500 slide evaluations, 1{,}741 Q\&A evaluations, 50 human ratings, plus the original baseline evaluations---an order-of-magnitude expansion of the empirical scope of the original submission.

\bigskip
\hrule
\bigskip

\reviewer{1}

\rcomment{The contribution lacks generalizability: the experiments are limited to a single downstream task, one generation model, and a relatively small dataset of 101 documents.}

\rresponse{I agree that single-model, single-task, single-dataset studies risk being unrepresentative. The revised paper addresses each axis directly:
\begin{itemize}
    \item Three generation models: GPT-4o, Claude Sonnet~4, and Gemini~2.0~Flash, each tested on every (document, method) pair. Method rankings are consistent across all three (see ``Cross-Generator Consistency'' subsection in Results).
    \item Two downstream tasks: slide generation (subjective quality) and factual Q\&A (objective correctness, see ``Question Answering: Truncation Dominates'' subsection in Results). The two tasks reveal qualitatively different chunking sensitivities, which is itself a key new contribution.
    \item Expanded corpus of 218 documents across more than 20 domains, including 74 non-research documents to test generalizability beyond the academic-paper genre (see ``Dataset'' subsection in Methodology).
\end{itemize}}

\rcomment{Recommendation 1: Strengthen the evaluation methodology with human evaluation or by validating LLM-as-judge against human judgments.}

\rresponse{A human evaluation study on 50 randomly selected document pairs has been added (see the new ``Human Evaluation Validates LLM Judges'' subsection in Results). For each pair, an expert annotator was shown two anonymized slide decks side-by-side---one from position-aware chunking and one from a baseline---and asked to choose which deck was better on six dimensions plus overall preference. The answer key was sealed until ratings were complete. Results confirm the LLM judges' findings: position-aware chunking is preferred in only 38\% of overall comparisons. This rules out single-judge bias as the explanation for position-aware chunking's lower scores.}

\rcomment{Recommendation 2: Expand the experimental scope by evaluating on additional downstream tasks (long-document summarization, question answering, RAG).}

\rresponse{The revision adds factual question answering as a second downstream task (see ``Question Answering'' subsection in Methodology and ``Question Answering: Truncation Dominates'' subsection in Results). For each source document, 10 factual questions with ground-truth answers are generated by Gemini~2.0~Flash from the full PDF before any chunking is applied. Each chunked text is then presented to a generation model, which must answer the 10 questions using only the chunked content. Per-method accuracy is computed objectively. Q\&A results reveal dramatically different chunking sensitivities than slides: an 18 percentage-point gap between best and worst, versus 1.7\% on slides. RAG over multi-document collections is acknowledged as a complementary direction in the ``Limitations and Scope'' subsection, but is qualitatively different from this setting because chunks in RAG are selected per query at inference time rather than per document offline.}

\rcomment{Recommendation 3: Test multiple models. The study currently uses only GPT-4o for generation.}

\rresponse{The revision uses three production-grade LLMs from three independent providers as generators: GPT-4o (OpenAI), Claude Sonnet~4 (Anthropic), and Gemini~2.0~Flash (Google), as described in the ``Generation Models'' subsection. All three reproduce the same method ranking (see ``Cross-Generator Consistency'' subsection in Results), confirming that the findings are not model-specific.}

\rcomment{Recommendation 4: Conduct ablation studies on parameters such as chunk budgets, chunk sizes, and prompt structures.}

\rresponse{Two ablation/sensitivity studies have been added, and prompt structure is explicitly addressed:
\begin{itemize}
    \item Chunk budget ablation (new ``Ablation: Chunk Budget Sensitivity'' subsection in Results): All four methods were re-run with budgets of 1{,}000, 2{,}000, 3{,}000, and 5{,}000 words on all 218 documents. Method rankings are stable across budgets.
    \item W-curve parameter sensitivity (new ``W-Curve Parameter Sensitivity'' subsection in Results): Position-aware chunking was re-run with 11 alternative W-curve configurations varying intro weight, conclusion weight, results-region peak, and floor (\emph{baseline}, \emph{intro\_low}, \emph{intro\_high}, \emph{concl\_low}, \emph{concl\_high}, \emph{res\_low}, \emph{res\_high}, \emph{floor\_zero}, \emph{floor\_high}, \emph{flat}, \emph{u\_curve}). Output size variance is under 2\% across configurations, demonstrating robustness to parameter perturbation.
    \item Prompt structure: All generation calls used the same prompt template across methods, generators, and tasks (held constant by design to isolate the chunking variable). This is acknowledged in the ``Limitations and Scope'' subsection as a design choice that leaves prompt-aware chunking as a complementary research direction.
\end{itemize}}

\rcomment{Recommendation 5: Provide deeper result analysis, particularly the consistent loss of statistical information.}

\rresponse{The new ``Statistics Retention: A Universal Weakness'' subsection in Results, together with the discussion, expands on this point. The universally poor statistics retention (approximately 2.6/5 across all judges and all generators) is attributed to a generation-side rather than chunking-side limitation: LLMs systematically prefer to summarize trends rather than cite specific numbers, regardless of whether the input chunk contains them. This is supported by the observation that all four chunking methods show similar statistics retention scores despite very different content selection strategies---if the limitation were on the chunking side, position-aware chunking (which explicitly rewards numeric content) would retain more statistics. This is discussed in the ``Why Does Truncation Win'' subsection, and prompt-side and post-processing techniques are identified as the relevant future work in the Conclusion.}

\rcomment{Recommendation 6: Clarify the limitations and scope of the findings, framing the conclusions more cautiously as task- and dataset-specific.}

\rresponse{The conclusions are now framed cautiously and explicitly tied to scope. The expanded ``Limitations and Scope'' subsection lists eight limitations (task coverage, prompt structure, domain coverage, question generation source, single human rater, language, static budgets, fixed evaluation rubric) and ends with: ``These limitations frame the findings as task-specific and dataset-specific rather than universal claims.'' The practitioner guide in the ``When to Use What'' subsection explicitly conditions recommendations on document length and downstream task type rather than offering a single universal recommendation. The conclusion likewise frames the takeaways as task-aware rather than universal, and the abstract uses phrases like ``has a clear niche'' to avoid overclaiming.}

\rcomment{Recommendation 7: Provide deeper theoretical reflection on why the observed results occur. For example, if truncation works well, is this because research papers are inherently front-loaded, because slide generation is a high-level summarization task, or because the evaluation method is insensitive to deeper content differences? Likewise, if position-aware chunking underperforms, does this suggest positional heuristics are fundamentally weak, or only that the current formulation is too simplistic?}

\rresponse{This is addressed in three places:
\begin{enumerate}
    \item A new ``Theoretical Framework'' section formalizes chunking as a budgeted utility-maximization problem with two components: information density along the document and LLM attention weight along the context. The W-shaped position curve is derived from the combination of these factors. Three conditions under which position-aware chunking should outperform truncation are stated explicitly: long documents, rich middle content, and LLMs that integrate discontinuous chunks. These theoretical predictions are then empirically tested in the ``Position-Aware Chunking Has a Niche'' subsection, where position-aware chunking's empirical sweet spot of 5{,}000 to 10{,}000-word documents is consistent with the theory.
    \item A new ``Why Does Truncation Win? Three Hypotheses'' subsection in the Discussion addresses the reviewer's three candidate explanations directly. The empirical results confirm Hypothesis~1 (front-loading), confirm Hypothesis~2 (slide generation is a high-level summarization task that does not stress fine-grained recall), and \emph{reject} Hypothesis~3 (evaluation insensitivity)---ruling it out using cross-judge agreement, human evaluation alignment, and the 18-percentage-point Q\&A gap that the same evaluation infrastructure produces.
    \item A new ``Why Does Position-Aware Chunking Underperform?'' subsection in the Discussion addresses the reviewer's question of whether position-aware chunking's underperformance reflects fundamental weakness of positional heuristics or only simplicity of the current formulation. The argument (supported by the medium-document niche and the W-curve sensitivity analysis) is that positional heuristics are not fundamentally weak, but the current formulation is overly simplistic in three identifiable ways: (1) document-structure assumption (no genre awareness), (2) discontinuity penalty (no coherence-aware selection), and (3) static budget allocation. Each is mapped to a future-work direction.
\end{enumerate}}

\bigskip
\hrule
\bigskip

\reviewer{2}

\rcomment{Recommendation 1: The W-shaped scoring function uses specific positional bounds (0.38--0.62) and quality signal weights without justification. Please either explain how these values were determined or include a brief sensitivity analysis.}

\rresponse{Both justification and sensitivity analysis are provided:
\begin{enumerate}
    \item Theoretical derivation (in the new ``Theoretical Framework'' section): The W-curve is derived from the combination of information density (which peaks in introduction, results, and conclusion regions for academic papers) and LLM attention weighting (which peaks at boundaries per Liu \emph{et al.} 2024).
    \item Sensitivity analysis (new ``W-Curve Parameter Sensitivity'' subsection in Results): Position-aware chunking was re-run with 11 alternative parameter configurations covering perturbations to every weight in the W-curve. Output size variance across all configurations is under 2\%, demonstrating that the qualitative shape (boundary emphasis with optional middle bump) matters more than the specific numeric values.
    \item Empirical validation (new ``Position-Aware Chunking Has a Niche'' subsection): The W-curve's predicted niche (medium-length documents with rich middle content) is empirically confirmed---position-aware chunking outperforms all baselines on slides for 5{,}000 to 10{,}000-word documents.
\end{enumerate}}

\rcomment{Recommendation 2: A small human evaluation on a subset (20--30 output pairs), especially for the Completeness and Balance dimensions, would strengthen the paper's main claims.}

\rresponse{Human evaluation was conducted on 50 document pairs (exceeding the suggested 20--30) in the new ``Human Evaluation Validates LLM Judges'' subsection. Ratings cover all six evaluation dimensions including Completeness and Coverage Balance. Results confirm the LLM judges' rankings.}

\rcomment{Recommendation 3: A paragraph discussing whether the results would be expected to hold for other generation tasks (summaries, Q\&A, report generation) and what conditions might change the conclusion would help readers assess generalizability.}

\rresponse{The revision adds an entire downstream task (Q\&A) rather than just a paragraph, providing direct empirical evidence on generalizability across tasks. The new ``Question Answering: Truncation Dominates'' subsection in Results shows that the chunking-method gap is dramatically larger on Q\&A (18 percentage points) than on slides (1.7\%). The new ``Task-Dependent Chunking'' subsection in the Discussion explicitly addresses when chunking methods matter and when they do not. RAG, report generation, and dialogue are acknowledged in the ``Limitations and Scope'' subsection as future work directions where conclusions might change.}

\rcomment{Recommendation 4: The 3{,}700$\times$ speed difference between semantic chunking and truncation is a strong practical argument against semantic chunking, but this implication could be stated more directly in the discussion and abstract.}

\rresponse{The speed argument has been made prominent in four places:
\begin{enumerate}
    \item Abstract: The closing sentence now explicitly states the measured 3{,}300$\times$ latency advantage of truncation over semantic chunking and frames it as a strong practical case for truncation as the default.
    \item New ``Computational Cost (Latency)'' subsection in Results: A dedicated subsection with a latency table reports per-document latency for all four methods, translates this into pipeline-scale impact (3.7 hours versus 4 seconds for 10{,}000 documents), and explicitly states that the latency penalty makes semantic chunking ``a poor default choice for throughput-sensitive applications.''
    \item New ``Practical Implications: The Speed-Quality Trade-off'' subsection in the Discussion: A new discussion subsection ties the latency results to the quality results and frames position-aware chunking as a much better complexity/quality trade-off than semantic chunking (approximately 74$\times$ faster than semantic chunking with comparable or better quality).
    \item Conclusion: A new paragraph summarizes the speed-quality trade-off as the second of the two practical takeaways of the paper.
\end{enumerate}}

\rcomment{Recommendation 5: A few implementation details needed for reproducibility are missing: the exact evaluation prompt given to Gemini, the chunk deduplication logic for the fixed-size approach, and the complete position-aware chunking selection algorithm.}

\rresponse{Each of the three specific items is now addressed:
\begin{enumerate}
    \item Exact evaluation prompt: A condensed excerpt of the multi-judge evaluation prompt now appears as a typeset block in the ``Multi-Judge LLM Evaluation for Slides'' subsection of Methodology, showing the six rated dimensions, JSON output schema, and document and slide-deck placeholders. The full prompt with formatting is available in the public repository linked in the manuscript.
    \item Chunk deduplication logic for fixed-size: A clarification has been added in the ``Fixed-Size First-Last'' description noting that fixed-size first-last does not employ a deduplication step because the first 1{,}000 and last 1{,}000 words are non-overlapping by construction (the dataset's minimum document length of 2{,}985 words ensures this). It is also clarified that semantic chunking does not need explicit deduplication because breakpoint detection ensures that adjacent chunks discuss different topics.
    \item Complete position-aware chunking algorithm: Algorithm~1 in the Methodology section provides the full pseudocode of the position-aware chunking selection procedure, including windowing, scoring, redundancy gate (Jaccard $\tau = 0.6$), coverage-ensuring step, and position-based reordering.
\end{enumerate}
In addition, all code, prompts, and corpus metadata are available at the GitHub repository linked in the Methodology section.}

\bigskip
\hrule
\bigskip

\section*{Closing}

I hope the revisions adequately address the reviewers' concerns. The cumulative empirical evidence---6{,}473 slide evaluations across three generators and three judges, 1{,}741 Q\&A evaluations with objective correctness, 50 human ratings, ablations on budgets and parameters, an expanded 218-document corpus spanning more than 20 domains, and a theoretical framework predicting and confirming position-aware chunking's empirical niche---substantially strengthens the original contribution.

The revised paper now reports a richer, more nuanced finding than the original: chunking method effects are task-dependent and length-dependent, and the paper provides practitioners with actionable guidance for when each method should be used. The original empirical observations (small overall differences on slides, statistics retention as a universal weakness) are preserved and reinforced, while the new evidence rules out the possibility that those observations were artifacts of a single model, task, judge, or dataset.

I thank the reviewers and editor for their time and constructive feedback.

\bigskip
\noindent\textit{Sincerely,}\\
Gourav Singla\\
Independent Researcher

\end{document}
